% !TeX root = ../DoABook.tex


\chapter{طرح‌های پیش‌پردازش و تخمین مرتبه مدل}
\label{chap:4}

\section{مقدمه}
\label{sec:4-1}
الگوریتم‌های تخمین \lr{DOA} توصیف‌شده در فصل قبل، فرض می‌کنند سیگنال‌های تابیده ناهمدوس\LTRfootnote{Noncoherent signals} هستند. با این حال، اگر سیگنال‌های تابیده شدیداً همبسته یا نسبت به یکدیگر همدوس\LTRfootnote{Coherent} باشند، اکثر الگوریتم‌ها در ارائه تخمین‌های قابل‌اطمینان \lr{DOA} شکست خواهند خورد؛ زیرا ماتریس کوواریانس داده $\bm{R}_{xx}$ دریافتی توسط آرایه، تکین\LTRfootnote{Singular} یا بدشرط\LTRfootnote{Ill-conditioned} می‌شود. این وضعیت در سناریوهای چندمسیره\LTRfootnote{Multipath} غیرمعمول نیست. در یک محیط غنی از چندمسیرگی، سیگنال‌های تابیده بر آرایه اغلب نسخه‌های تأخیریافته و مقیاس‌شده‌ای از یکدیگر هستند و بنابراین شدیداً همبسته یا همدوس می‌باشند. برای رفع این مشکل، ماتریس کوواریانس داده پیش از «تغذیه» به الگوریتم‌های تخمین \lr{DOA}، پردازش می‌شود. این تکنیک‌ها، روش‌های پیش‌پردازش نامیده می‌شوند و نقشی حیاتی، اگر نگوییم اجباری، در تخمین \lr{DOA} ایفا می‌کنند. دو طرح پیش‌پردازش شناخته‌شده و تثبیت‌شده، یعنی هموارسازی مکانی\LTRfootnote{Spatial smoothing} و میانگین‌گیری رفت‌وبرگشت\LTRfootnote{Forward-backward averaging}، در این فصل مورد بحث قرار خواهند گرفت.

علاوه بر مسئله همدوسی سیگنال، تاکنون فرض شده است که تعداد سیگنال‌های تابیده بر آرایه مشخص است. در واقعیت اما، این تعداد مشخص نیست و باید از روی داده‌های دریافتی تخمین زده شود. الگوریتم‌های تخمین \lr{DOA} کاملاً به فرضِ آگاهی از تعداد سیگنال‌های تابیده وابسته هستند. از این رو، تخمین این تعداد نقش کلیدی دیگری را در فرآیند تخمین \lr{DOA} ایفا می‌کند و پژوهش در رده مجزایی از تکنیک‌ها به نام تکنیک‌های تخمین مرتبه مدل\LTRfootnote{Model order estimation} را می‌طلبد. تخمین‌گرهای مرتبه مدل یا مرتبه منبع، الگوریتم‌هایی اجباری هستند که باید پیش از اجرای الگوریتم‌های تخمین \lr{DOA} اجرا شوند. این فصل مقدمه کوتاهی بر برخی از این تکنیک‌ها ارائه می‌دهد.

\section{طرح‌های پیش‌پردازش}
\label{sec:4-2}
الگوریتم‌های تخمین \lr{DOA} مبتنی بر زیرفضا، کاملاً بر شرط تمام‌رتبه\LTRfootnote{Full rank} بودن ماتریس کوواریانس داده $\bm{R}_{xx}$ که در رابطه \eqref{eq:4-1} آمده، استوار هستند. این رابطه در اینجا برای دسترسی سریع دوباره آورده شده است:
\begin{equation}
	\bm{R}_{xx} = E\{\bm{x}(t)\bm{x}^H(t)\} = \bm{A}\bm{R}_{ss}\bm{A}^H + \sigma_N^2 \bm{I}_M
	\label{eq:4-1}
\end{equation}
وضعیت $\bm{R}_{xx}$ به وضعیت ماتریس کوواریانس سیگنال $\bm{R}_{ss}$ بستگی دارد. زمانی که $d$ جبهه موج سیگنالِ تابیده با هم همبسته نباشند، ماتریس کوواریانس سیگنال $\bm{R}_{ss}$ دارای رتبه کامل $d$ است؛ در این حالت ماتریس قطری و غیرتکین است. اما هنگامی که سیگنال‌ها تا حدودی همبسته باشند، یا زمانی که برخی سیگنال‌ها کاملاً همبسته (همدوس) باشند، $\bm{R}_{ss}$ غیرقطری و تکین می‌شود. برای مثال، اگر یکی از جبهه‌های موجِ سیگنالِ تابیده شدیداً با دیگری همبسته یا با آن همدوس باشد، بزرگی ضریب همبستگی متقابل آن‌ها برابر با ۱ شده و رتبه $\bm{R}_{ss}$ به $d-1$ کاهش می‌یابد. بنابراین، شرط تمام‌رتبه بودن در الگوریتم‌های \lr{DOA} که پیش‌تر توصیف شد، دیگر برقرار نخواهد بود (یعنی $\text{rank}\{\bm{A}\bm{R}_{ss}\bm{A}^H\} = d - 1$)؛ در نتیجه، الگوریتم \lr{DOA} شکست خواهد خورد. به‌طور کلی، اگر $P$ جبهه موج همدوس وجود داشته باشد، رتبه $\bm{R}_{ss}$ برابر با $M-P$ خواهد بود. در چنین حالتی، $\bm{R}_{xx}$ دارای $d$ مقدار ویژه مثبت بزرگ نخواهد بود (برخلاف حالتی که سیگنال‌ها ناهمبسته هستند و $\bm{R}_{ss}$ قطری است). این امر به این معنی است که ماتریس $\bm{R}_{ss}$ احتمالاً کمتر از $d$ مقدار ویژه مثبت بزرگ داشته باشد که موجب بروز خطا در فرآیندهای تخمین \lr{DOA} می‌گردد.



برای مثال، فرض کنید دو مورد از سیگنال‌های تابیده، همدوس باشند؛ مثلاً دو سیگنال اول همدوس هستند، یعنی $s_2(t) = \alpha s_1(t)$ که در آن $\alpha$ یک اسکالر مختلط است که رابطه بهره و فاز بین دو سیگنال همدوس را توصیف می‌کند. در این صورت $\bm{s}(t)$ بردار سیگنالی است که به‌صورت زیر داده می‌شود:
\begin{equation}
	\bm{s}(t) = [s_1(t), \alpha s_1(t), s_3(t), \dots, s_d(t)]^T
	\label{eq:4-2}
\end{equation}
و ماتریس هدایت آرایه به‌صورت زیر خواهد بود:
\begin{equation}
	\bm{A} = [\bm{a}(\theta_1), \bm{a}(\theta_2), \bm{a}(\theta_3), \dots, \bm{a}(\theta_d)]
	\label{eq:4-3}
\end{equation}
با این سیگنال‌های تغییریافته، ماتریس کوواریانس سیگنال $\bm{R}_{ss} = E\{\bm{s}(t)\bm{s}^H(t)\}$ اکنون یک ماتریس $d \times d$ تکین است و دیگر غیرتکین نیست. دو خطای زیر در این حالت قابل مشاهده است:
\begin{enumerate}
	\item مرتبه تکرار\LTRfootnote{Multiplicity} کوچک‌ترین مقادیر ویژه $\bm{R}_{xx}$ اکنون برابر با $k = M - d + 1$ است و نه $k = M - d$. این امر منجر به بروز خطا در تخمین تعداد سیگنال‌ها به‌صورت $d - 1 (= M - k)$ به‌جای $d$ می‌شود.
	\item به دلیل کاهش تعداد مقادیر ویژه بزرگِ قابل‌تفکیک در $\bm{R}_{xx}$، تنها $\{\theta_3, \dots, \theta_d\}$ قابل تفکیک هستند.
\end{enumerate}
برای مقابله با مسئله دریافت سیگنال همدوس، می‌توان از طرح‌های پیش‌پردازشی مانند میانگین‌گیری رفت‌وبرگشت یا هموارسازی مکانی استفاده کرد؛ این طرح‌ها تضمین می‌کنند که $\bm{R}_{ss}$ (پس از هموارسازی) حتی زمانی که تمام سیگنال‌های دریافتی همدوس هستند، تمام‌رتبه و غیرتکین باشد. به عبارت دیگر، این طرح‌ها اساساً برای ناهمبسته کردن\LTRfootnote{Decorrelate} سیگنال‌ها پیش از تخمین \lr{DOA} آن‌ها به کار می‌روند \cite{4-1, 4-2}.

\subsection{میانگین‌گیری رفت‌وبرگشت}
\label{subsec:4-2-1}
میانگین‌گیری رفت‌وبرگشت\LTRfootnote{Forward-backward (FB) averaging} یک طرح پیش‌پردازش محبوب است که در بسیاری از کاربردهای پردازش سیگنال از جمله جهت‌یابی رادیویی\LTRfootnote{Radio direction finding} به کار گرفته می‌شود. میانگین‌گیری رفت‌وبرگشت در صورتی قابل پیاده‌سازی است که آرایه‌ها ماهیتاً مرکز-متقارن باشند. عملیات اصلی میانگین‌گیری رفت‌وبرگشت بر این واقعیت استوار است که بردارهای هدایت \lr{ULA} حتی در صورت معکوس شدن ترتیب عناصر و مزدوج مختلط شدن آن‌ها، یکسان باقی می‌مانند.

فرض کنید $\bm{\Pi}_M$ یک ماتریس تعویض $M \times M$ باشد. در این صورت برای \lr{ULA} می‌توان نشان داد که \cite{4-1}:
\begin{equation}
	\bm{\Pi}_M \bar{\bm{a}}(\theta_i) = e^{-j(M-1)\mu_i} \bm{a}(\theta_i)
	\label{eq:4-4}
\end{equation}
سپس می‌توان یک ماتریس کوواریانس داده‌ی بازگشتی\LTRfootnote{Backward data covariance matrix} $\bm{R}_{back}$ را از ماتریس کوواریانس داده‌ی واقعی $\bm{R}_{xx}$ به‌صورت زیر تشکیل داد:
\begin{equation}
	\bm{R}_{back} = \bm{\Pi}_M \bar{\bm{R}}_{xx} \bm{\Pi}_M
	\label{eq:4-5}
\end{equation}

مورد ساده‌ای از دو منبع همدوس که بر یک \lr{ULA} می‌تابند را در نظر بگیرید. ماتریس کوواریانس داده‌ی میانگین‌گیری‌شده‌ی رفت‌وبرگشت از طریق میانگین‌گیری از ماتریس کوواریانس واقعی (رفت) و ماتریس کوواریانس برگشت به‌دست می‌آید، یعنی:
\begin{align}
	\bm{R}_{xx}^{fb} &= \frac{1}{2} (\bm{R}_{xx} + \bm{R}_{back}) = \frac{1}{2} (\bm{R}_{xx} + \bm{\Pi}_M \bar{\bm{R}}_{xx} \bm{\Pi}_M) \nonumber \\
	&= \frac{1}{2} E\{\bm{x}(t)\bm{x}^H(t)\} + \frac{1}{2} \bm{\Pi}_M \overline{E\{\bm{x}(t)\bm{x}^H(t)\}} \bm{\Pi}_M \nonumber \\
	&= \bm{A} \left[ \frac{1}{2} (\bm{R}_{ss} + \bm{\Delta} \bar{\bm{R}}_{ss} \bm{\Delta}^H) \right] \bm{A}^H + \sigma_N^2 \bm{I}_M
	\label{eq:4-6}
\end{align}
که در آن $\bm{\Delta} \in \mathbb{C}^{d \times d}$ یک ماتریس قطری یکانی است \cite{4-3}. با مقایسه رابطه \eqref{eq:4-6} با \eqref{eq:4-1}، ماتریس کوواریانس سیگنال جدید به‌صورت زیر داده می‌شود:
\begin{equation}
	\bm{R}_{ss}^{fb} = \frac{1}{2} (\bm{R}_{ss} + \bm{\Delta} \bar{\bm{R}}_{ss} \bm{\Delta}^H)
	\label{eq:4-7}
\end{equation}
این ماتریس کوواریانس سیگنال رفت‌وبرگشت جدید، حتی اگر دو سیگنال همدوس باشند، دارای رتبه $d$ است؛ این امر جداسازی دو سیگنال همدوس یا شدیداً همبسته را امکان‌پذیر می‌سازد. بنابراین، نسخه‌ی میانگین‌گیری‌شده‌ی رفت‌وبرگشتِ هر الگوریتمِ مبتنی بر کوواریانس را می‌توان با جایگزینی $\bm{R}_{xx}$ با $\bm{R}_{xx}^{fb}$ در یک الگوریتم \lr{DOA} به‌دست آورد. این تکنیک گاهی حتی در محیط‌های با سیگنال‌های ناهمبسته نیز برای دستیابی به تخمین‌های قابل‌اطمینان‌تر استفاده می‌شود.






























